\section{Characterization Results}
\subsection{Co-run loss under original wiring}
\paperfigure{fig01-throughput}{Victim throughput and elapsed-time contrasts from the fixed repaired baseline. DS0 is the primary configuration. The A1 co-run term and the additional higher-working-set term are descriptive contrasts, not a measured partition of PTW cost.\claim{R-ELAPSED}\claim{M-CONTROLS}}{fig:throughput}
Figure~\ref{fig:throughput} establishes the main result without inserted data-structure-path delay. Solo measures the fixed victim alone. A1 adds a co-running neighbor while both streams retain zero HPM miss counts.\claim{R-MISS} The Solo--A1 victim-throughput decrease therefore cannot be attributed to measured misses. The comparison is essential even though it is a conventional control rather than a methodological invention.

Increasing the neighbor's configured working set creates additional loss. Under DS0, A8 and A64 add \DSZeroAEightElapsedExtra\% and \DSZeroASixtyFourElapsedExtra\% to the A1 median victim hardware elapsed time.\claim{R-ELAPSED} The denominator is the matching A1 result, not Solo. Hardware throughput is based on the same fixed useful victim payload in each trial; a percentage increase in elapsed time and a percentage decrease in throughput have different denominators.\claim{M-STAT}

Victim slowdown must be distinguished from aggregate progress. Using completed aggressor write responses as a common-window payload proxy, A1's aggregate progress is \DSZeroAOneAggregatePercent\% of the matched Solo throughput. A8 and A64 instead show aggregate proxy losses of \DSZeroAEightAggregateLoss\% and \DSZeroASixtyFourAggregateLoss\%, respectively.\claim{R-AGGREGATE} The proxy counts useful COPY payload once and does not measure exact physical memory bandwidth. The underlying record does not capture exact same-window write-data bytes; its documented transaction-boundary error remains attached to the estimate.\claim{M-AGGREGATE} Thus a victim receiving roughly half of concurrent progress is different from a loss of half of system throughput.

These measurements support a bounded claim about this implementation and offered load: a neighbor with a larger translation working set can further delay a fixed victim beyond the zero-miss co-run loss. The contrast also changes the aggressor's address stream and service timing. It cannot isolate the incremental cycles spent walking page tables, the associated memory traffic, and internal serialization into separately additive costs.\claim{M-CONTROLS}

The HPM boundary matters. The DS0 A8, A16, and A64 conditions retain victim miss medians of two, while A1 retains zero on both sides.\claim{R-MISS} These samples are not discarded. Because the counters use a software snapshot interval, the observations neither establish a zero-miss victim for those DS0 conditions nor identify where the small counts fall relative to the hardware timer. The direct request-level evidence below uses a separate observation of the victim's first lookup.

\subsection{Where request lifetime grows}
\paperfigure{fig02-path-miss}{Segmented victim read/write lifetimes and separately measured HPM misses. Each path value is a cross-trial median of a per-trial request mean. Stages and channels overlap; the panel is not an additive time breakdown.\claim{M-PATH}\claim{M-STAT}}{fig:path}
Figure~\ref{fig:path} distinguishes the shared data-path effect from internal waiting. In the historical path experiment, Solo--A1 leaves the internal mean unchanged while increasing later request lifetime. A1--A8 instead increases internal passage markedly; the translated-address output does not acquire extra observed address-handshake backpressure.\claim{R-PATH} The new baseline reproduces the same qualitative separation, with a stronger internal increase under DS50.\claim{R-PATH-NEW}

The observer localizes a problem but does not itself establish its complete cause. Internal passage includes queueing before lookup and service, the translation operation, and exit-related waiting. The write completion stage also includes write-data availability and the return response.\claim{M-PATH} A delayed translated read can affect later write-data timing, while memory traffic can compete independently. We therefore retain the read and write panels and do not label every downstream increment as a DDR service increase.

The path mean is not a tail statistic. A larger mean may reflect many moderately delayed requests or a smaller set of long waits. The saved sums, counts, and finite maxima do not determine the distribution between them.\claim{M-STAT} This distinction constrains both the current conclusions and the pending latency-distribution experiment.

\subsection{A named response extends a neighboring hit's wait}
The historical E3 experiment addresses the local dependency left open by aggregate path counters. It selects one aggressor stage-one PTW response while a victim read is waiting, then compares delay-enabled and delay-disabled runs with recorded request identities and prefix events. The victim subsequently reports a first-lookup hit.\claim{R-E3}

Across the four formal complementary cold pairs, the added aggressor remaining service and victim remaining wait match pair by pair, ranging from \EThreeDeltaMin{} to \EThreeDeltaMax{} cycles. The victim's own observed translation-service interval remains \EThreeOwnService{} cycles.\claim{R-E3} This is evidence that the extended aggressor response delays the waiting victim before the victim's own hit service in the tested controller. Source inspection identifies the shared serial control state and read-address queue as the corresponding service dependency.\claim{R-E3-RTL}

This causal claim remains local. Only one complementary pair has matching recorded pre-trigger timing; the other pairs are retained with their differences. The hot controls do not trigger the delay and are negative controls for triggering, rather than equal-exposure tests.\claim{R-E3-LIMIT} The diagnostic uses a different transfer size and an historical image. Its cycle deltas do not explain the complete sustained-batch slowdown, and the experiment does not imply that the RISC-V architecture or Advanced eXtensible Interface (AXI) protocol requires this organization.

\subsection{Sensitivity to data-structure-path delay}
\paperfigure{fig03-ds-sweep}{Data-structure-path delay sensitivity. Completed DS0 and DS50 baseline points are distinguished from any pending DS25/DS100 measurements. A missing point is not interpolated into a measured curve.\claim{G-DS-SWEEP}\claim{M-DS}}{fig:ds}
The additional delay configuration magnifies the higher-working-set contrast. A8 and A64 increase median hardware elapsed time over DS50's A1 by \DSFiftyAEightElapsedExtra\% and \DSFiftyASixtyFourElapsedExtra\%, respectively, compared with the smaller DS0 increments.\claim{R-ELAPSED} Figure~\ref{fig:ds} reports configurations separately, preserving their actual collection order and boots.

The sensitivity is informative because translation-related service depends on the latency of the memory system and control path around it. It is not a clean PTW-only intervention: data-structure requests include other accesses, and altered schedules can change both aggressor and victim miss observations.\claim{M-DS}\claim{R-MISS} The experiment therefore supports sensitivity to the configured path delay; it does not estimate a universal slowdown per PTW cycle or demonstrate a page-table-priority policy.

\subsection{Coverage still required for generalization}
\paperfigure{fig04-device-count}{Device-count scaling slot. Four DMA engines are instantiated, but no completed two-to-three-to-four active-device matrix is available for this draft.\claim{G-COVERAGE}}{fig:devices}
\paperfigure{fig05-workload-size}{Workload-shape and transfer-size slot. Sequential, random, and periodic-burst comparisons require a matched baseline; no illustrative values substitute for measurements.\claim{G-COVERAGE}}{fig:shapes}
\paperfigure{fig06-page-size}{Page-size comparison slot. A valid comparison must verify the actual page-table leaf level and alignment; buffer allocation or address stride alone does not establish a larger translation page.\claim{G-COVERAGE}}{fig:pages}
\paperfigure{fig07-latency-cdf}{Request-latency distribution slot. Existing means and limited observed maxima cannot reconstruct a CDF or p99.\claim{G-COVERAGE}\claim{M-STAT}}{fig:cdf}
Figures~\ref{fig:devices}--\ref{fig:cdf} preserve the planned tests whose results are not present. The active-device count, the aggressor's individual and total offered pressure, and the completion of every stream must remain distinguishable. Workload-shape and transfer-size studies must retain matched work and expose the supply pattern. Page-size studies must verify the installed translation leaf, not infer it from a DMA buffer. Request-latency percentiles require bounded distribution capture with overflow or loss accounting.\claim{G-COVERAGE}

These gaps limit comparison with other platforms. Bermuda includes transfer-size-dependent results, so a batch of larger copies on our platform cannot be treated as a direct contradiction of its curve~\cite{bermuda2025}.\claim{L-BERMUDA} The unresolved question is which service and supply conditions preserve or reduce interference on this implementation. The working draft makes no device-scaling, page-size, or latency-tail finding before those measurements exist.

\subsection{Historical mitigation costs}
\paperfigure{fig08-mitigation-history}{Historical mitigation evidence. Victim protection and aggressor progress retain their own common-window denominators. Separate version and collection groups are not pooled into one new-baseline matrix.\claim{R-MITIGATION}}{fig:mitigation}
The historical rate-control experiment improves victim progress by reducing the aggressor's issue pressure, with a corresponding cost to aggressor progress in the common window.\claim{R-MITIGATION} Aggressor work completed after the victim window is not divided by that window's duration and relabeled concurrent throughput. This accounting prevents protection of one stream from appearing free because the other's tail is ignored.

A distinct composition experiment compares partitioning (P), bypass-related service (N), and their combination (PN). After isolating replacement history, PN remains \PNPenaltyMin\%--\PNPenaltyMax\% slower than P across three descriptive paired comparisons even though their aggressor miss counts match.\claim{R-PN} Equal miss counts do not establish equal walk addresses, memory-service schedules, or request timing. The full cause of this board-level difference remains unresolved. A reachable contention path in a behavioral memory model is not proof that it explains the measured board difference.\claim{R-PN-LIMIT}

These are configuration-specific tradeoffs and a preserved negative result. They do not establish that partitioning or rate control is a new technique, that composing optimizations is generally harmful, or that the historical results already form a matched mitigation matrix on the repaired baseline.\claim{R-MITIGATION-LIMIT} The latter remains a separate measurement gap.

\subsection{Capacity replacement is a separate effect}

The historical E2 experiments observe residency before a controlled access sequence, removal under one boundary layout, and a miss on the first later reuse of the victim translation.\claim{R-CAPACITY} This supports a capacity-replacement chain under those layouts. It differs from E3, where a victim request that first hits waits for an aggressor's service. Cache occupancy and service serialization must therefore be considered separately even when both accompany a larger neighbor working set.

The tested W7/W8 layouts do not define a universal device capacity threshold. Source and destination translation keys, access order, replacement state, and direction are part of the experimental condition.\claim{R-CAPACITY-LIMIT} We retain this finite-layout evidence rather than use it to partition the new baseline's total slowdown.

\subsection{Capacity sensitivity: separate register-transfer-level simulation}
\paperfigure{fig09-capacity}{IOTLB-capacity sensitivity in a separate real register-transfer-level (RTL) simulation. Its software-register-fed workload and behavioral memory service differ from autonomous FPGA batch supply. No repeated-sample confidence interval or FPGA time conversion is implied.\claim{R-CAPACITY-SIM}}{fig:capacity}
Figure~\ref{fig:capacity} reports the completed capacity sweep from the separate real-RTL simulation package. The simulation uses ordinary memory-mapped I/O (MMIO) register issue and a single COPY task in flight rather than the board's autonomous batch fixture. The victim warm-up uses 64 copies plus two zero-walk checks, rather than the board's 256-copy warm-up; no CPU or Linux executes in this model. The memory model serves reads in arrival order and waits 16 cycles at a read-queue head, then returns one 8-byte beat per cycle. Its baseline-capacity calibration reproduces the ordering of Solo, A1, and A8, while retaining different absolute timing and memory service from the board.\claim{R-CAPACITY-SIM}

For A8, increasing capacity eliminates the aggressor walks at the intermediate and largest capacities, and victim elapsed time becomes equal to the matching A1 time. The A1--Solo difference remains. The larger A64 working set retains aggressor walks and additional delay at the largest evaluated capacity.\claim{R-CAPACITY-SIM} These results distinguish a capacity-sensitive increment from remaining co-run delay in this model; they do not separate that residual into serial-controller waiting and memory competition. The simulation cannot replace the historical board residency experiment or establish a capacity benefit for unmeasured FPGA configurations.

\subsection{Invalidation scope: simulation supplement}
\paperfigure{fig10-invalidation}{Invalidation-scope results from Verilator with behavioral DDR service. Time axes are simulation steps. These are neither Linux execution nor FPGA performance results.\claim{R-INVALIDATION}}{fig:invalidation}
The separate invalidation study runs real RTL DMA and command-queue logic against a behavioral DDR service model. It varies invalidation scope and domain while accounting for the issuing stream's complete work, final flush, and the neighbor's complete work separately.\claim{R-INVALIDATION} A narrow invalidation can eliminate the misses attributed to removed neighboring resident translations, yet increase the issuing stream's complete time relative to a wider invalidation.\claim{R-INVALIDATION}

Attribution uses the recorded resident-key invalidation and first subsequent miss, rather than counting every extra walk over a baseline as invalidation-induced.\claim{R-INVALIDATION-METHOD} The supplement tests an explicit scope/command tradeoff in the model. It does not supply a current-mainline Linux cost, a hardware time conversion, or an FPGA validation of the simulated conclusions.\claim{R-INVALIDATION-LIMIT}
